
\documentclass[12pt,a4paper]{article}

% Configuración del documento
\usepackage[utf8]{inputenc}
\usepackage[T1]{fontenc}
\usepackage[spanish]{babel}
\usepackage{geometry}
\usepackage{longtable}
\usepackage{array}
\usepackage{booktabs}

\geometry{margin=2.5cm}

\begin{document}

\begin{center}
    {\Large \textbf{Universidad Técnica Estatal de Quevedo}}\\[0.3cm]
    \textbf{Carrera de Software}\\
    \textbf{Ingeniería de Requisitos (ISR-401)}\\[0.5cm]

    {\large \textbf{FORM-TL-01}}\\
    Taller de identificación y clasificación de tipos de requisitos
\end{center}

\vspace{0.5cm}

\noindent
\textbf{Estudiante:} Luis Yepez\\
\textbf{Periodo académico:} 2026--2027\\
\textbf{Tipo de actividad:} Formativa\\
\textbf{Caso de estudio:} Plataforma de gestión de citas médicas

\section{Objetivo}

Identificar y clasificar los requisitos de una plataforma de
gestión de citas médicas, aplicando los fundamentos de
Ingeniería de Requisitos.

\section{Clasificación de requisitos}

En esta sección se analizarán los 20 enunciados
proporcionados en el caso de estudio.

\section{Análisis de requisitos de calidad}

Se identificarán cinco requisitos de calidad
y se propondrán redacciones verificables.

\section{Discusión de resultados}

Se registrarán las diferencias de clasificación
identificadas durante la comparación con un compañero.

\section{Referencias bibliográficas}

Las referencias utilizadas se incorporarán
siguiendo el formato IEEE.

\end{document}
